% from ega4-9-fr.tex
\oldpage[IV-3]{84}

et (8.5.8), appliqués suivant la méthode de (8.1.2, \emph{a)}), il existe un
sous-$\mathcal O_X$-Module cohérent $\mathcal G$ de $\mathcal F$ tel que
$\mathcal G_\eta=\mathcal H$~; si $Y=\operatorname{Supp}(\mathcal F)$,
$Y'=\operatorname{Supp}(\mathcal G)$, on a par suite
$Y_s=\operatorname{Supp}(\mathcal F_s)$, $Y'_s=\operatorname{Supp}(\mathcal G_s)$
pour tout $s\in S$ (\textbf{I}, 9.1.13) et en particulier $Z=Y_\eta$,
$Z'=Y'_\eta$~; comme $Z-Z'$ est dense dans $Z$ et $Z'\neq\varnothing$, il y a
un voisinage $U$ de $\eta$ tel que pour $s\in U$, $Y_s-Y'_s$ soit dense dans
$Y_s$, et $Y'_s\neq\varnothing$ (9.5.1 et 9.5.3)~; en considérant un point
générique d'une composante irréductible de $Y'_s$ et un voisinage assez petit de
ce point dans $X_s$, on déduit aussitôt de (3.1.3) que $\mathcal F_s$ admet des
cycles premiers associés immergés.

\phantomsection\label{IV.9.8.2-fr}%
\textbf{(9.8.2)} Soient $S$ un préschéma noethérien \emph{intègre}, de point
générique $\eta$, $f:X\to S$ un morphisme de type fini, $\mathcal F$ un
$\mathcal O_X$-Module cohérent. Considérons une décomposition irredondante
réduite $(\mathcal G_i)_{i\in I}$ de $\mathcal F_\eta$ (3.2.6)~; les
$\mathcal G_i$ sont donc des quotients de $\mathcal F_\eta$ et il y a un
homomorphisme injectif
$h:\mathcal F_\eta\to\bigoplus_i\mathcal G_i$~; en outre
$\operatorname{Ass}(\mathcal G_i)$ est réduit à un point $x_i$. Soit $Z_i$
l'adhérence de $\{x_i\}$ dans $X$, de sorte que
$(Z_i)_\eta=\operatorname{Supp}(\mathcal G_i)$. Désignons par $\mathcal J_i$
l'Idéal cohérent de $\mathcal O_X$ définissant le sous-préschéma fermé réduit de
$X$ d'espace sous-jacent $Z_i$ (sous-préschéma encore noté $Z_i$). En vertu de
(8.5.2) et (8.5.8), appliqués suivant la méthode de (8.1.2, \emph{a)}), on peut
(en restreignant au besoin $S$ à un voisinage de $\eta$) supposer qu'il existe
des quotients $\mathcal F_i$ de $\mathcal F$ ($i\in I$) tels que
$\mathcal G_i=(\mathcal F_i)_\eta$ pour tout $i$, et un homomorphisme
$g:\mathcal F\to\bigoplus_i\mathcal F_i$ tel que $h=g_\eta$. En outre
(\textbf{I}, 9.3.5), il existe un entier $m$ tel que
$(\mathcal J_i^m\mathcal F_i)_\eta=(\mathcal J_i)_\eta^m\mathcal G_i=0$ et en
restreignant de nouveau $S$, on peut donc supposer que l'on a aussi
$\mathcal J_i^m\mathcal F_i=0$ (8.5.2.5), donc que le support de $\mathcal F_i$
est contenu dans $Z_i$~; mais comme il est fermé et contient $x_i$, il est
nécessairement égal à $Z_i$.

\begin{proposition}[9.8.3]
\label{IV.9.8.3-fr}
Sous les conditions de (9.8.2), pour tout $s\in S$, et tout $i\in I$, désignons
par $x_{is\alpha}$ ($\alpha\in J_{s,i}$) les points maximaux de $(Z_i)_s$. Il
existe un voisinage $U$ de $\eta$ dans $S$ tel que, pour tout $s\in U$, les
$x_{is\alpha}$ (pour $i\in I$ et, pour chaque $i$, $\alpha\in J_{s,i}$) soient
deux à deux distincts et que $\operatorname{Ass}(\mathcal F_s)$ soit l'ensemble
des $x_{is\alpha}$ (autrement dit, les cycles premiers associés à $\mathcal F_s$
sont les composantes irréductibles des $(Z_i)_s$). En outre, on peut prendre
$U$ tel que, pour que l'adhérence de $\{x_{is\alpha}\}$ dans $X_s$ soit un cycle
premier associé maximal de $\mathcal F_s$, il faut et il suffit que $(Z_i)_\eta$
(adhérence de $\{x_i\}$ dans $X_\eta$) soit un cycle premier associé maximal de
$\mathcal F_\eta$.
\end{proposition}

Il résulte de (3.1.3, \emph{c')}) que pour chaque $i$, il existe un ouvert $W_i$
dans $X_\eta$ tel que $W_i\cap(Z_i)_\eta$ soit non vide, et un
$\mathcal O_{X_\eta}$-Module cohérent $\mathcal H_i$, de support
$W_i\cap(Z_i)_\eta$, tels qu'il y ait un homomorphisme injectif
$v_i:\mathcal H_i\to\mathcal F_\eta|W_i$. Soit $V_i$ un ouvert de $X$ tel que
$V_i\cap X_\eta=W_i$~; appliquant, comme dans (9.8.2), les résultats de (8.5.2)
et (8.5.8), on peut (en restreignant $S$) supposer qu'il existe un Module
cohérent $\mathcal F'_i$, de support $V_i\cap Z_i$ et un homomorphisme
$u_i:\mathcal F'_i\to\mathcal F|V_i$ tels que
$\mathcal H_i=(\mathcal F'_i)_\eta$ et $v_i=(u_i)_\eta$. Nous allons démontrer
qu'il y a un voisinage $U$ de $\eta$ tel que pour $s\in U$, les propriétés
suivantes soient vérifiées~:
\begin{enumerate}[label=\emph{(\roman*)}]
  \item Les $(\mathcal F_i)_s$ sont sans cycle premier associé immergé.
  \item L'homomorphisme
  $g_s:\mathcal F_s\to\bigoplus_i(\mathcal F_i)_s$ est injectif.
  \item $(Z_i)_s\cap(V_i)_s$ est dense dans $(Z_i)_s$,
  $(\mathcal F'_i)_s$ a pour support $(Z_i)_s\cap(V_i)_s$ et
  $(u_i)_s:(\mathcal F'_i)_s\to\mathcal F_s|(V_i)_s$ est injectif.
\end{enumerate}

